\section{Method}\label{method}

\subsection{Overall architecture}\label{overall-architecture}

Let \(I \in \mathbb{R}^{3\times H\times W}\) denote an input colonoscopy image. The proposed
network follows a U-shaped encoder-decoder configuration and consists of
five principal components: a ConvNeXt-Tiny encoder, a lightweight
residual bottleneck, a Multi-Scale Context module, a boundary-sensitive
decoder with refined skip connections, and a high-resolution Detail
Branch with gated fusion.

The encoder produces four hierarchical feature maps \{F1, F2, F3, F4\}
with channel dimensions \{96, 192, 384, 768\} and spatial resolutions
\{H/4×W/4, H/8×W/8, H/16×W/16, H/32×W/32\}. The deepest feature is
processed by the lightweight bottleneck and the Multi-Scale Context
module. The decoder then progressively increases spatial resolution. At
the first three decoding stages, the upsampled decoder feature is
concatenated with the corresponding encoder representation and processed
by BSEI. At one-half input resolution, the final decoder feature is
fused with the output of the Detail Branch. A refinement head converts
the fused representation into final segmentation logits, which are
bilinearly resized to the original image resolution. During training,
three auxiliary prediction heads produce additional full-resolution
logits from intermediate decoder stages.

The principal forward operations can be summarized as follows:

\emph{\{F1, F2, F3, F4\} = E(I),}

\emph{B = M\_MSC(M\_LB(F4)),}

\emph{D4 = M\_BSEI,4({[}U4(B), F3{]}),}

\emph{D3 = M\_BSEI,3({[}U3(D4), F2{]}),}

\emph{D2 = M\_BSEI,2({[}U2(D3), F1{]}),}

\emph{D1 = M\_BSEI,1(U1(D2)),}

\emph{P = U\_out(M\_ref(M\_GDF(D1, M\_detail(I)))).}

Here, {[}·,·{]} denotes channel-wise concatenation, Ui denotes the
corresponding decoder up-sampling operation, M\_LB is the lightweight
bottleneck, M\_MSC is the Multi-Scale Context module, M\_GDF is Gated
Detail Fusion, and P denotes the final segmentation logits.

\subsection{\texorpdfstring{ ConvNeXt-Tiny
encoder}{ ConvNeXt-Tiny encoder}}\label{convnext-tiny-encoder}

Before entering the encoder, the input image is normalized using the
ImageNet channel mean and standard deviation:

\emph{\(\hat{I}_c = (I_c - \mu_c) / \sigma_c,\)}

\emph{\(\mu = (0.485, 0.456, 0.406), \sigma = (0.229, 0.224, 0.225).\)}

The first encoder stage uses a 4×4 convolution with stride 4, followed
by Layer Normalization. The subsequent three stages begin with Layer
Normalization and a 2×2 stride-2 convolution. The numbers of ConvNeXt
blocks in the four stages are 3, 3, 9, and 3, respectively.

For an input feature X, each ConvNeXt block applies a 7×7 depthwise
convolution, converts the representation to channels-last format for
Layer Normalization and linear channel projection, expands the channel
dimension by a factor of four, applies GELU activation, restores the
original channel dimension, and adds the transformed representation
through a layer-scaled stochastic residual connection.

\emph{Z = DWConv7×7(X),}

\emph{V = W2(GELU(W1(LN(Z)))),}

\emph{\(Y = X + \operatorname{DropPath}(\gamma \odot V).\)}

Spatial dropout is applied to the outputs of the second, third, and
fourth encoder stages. This regularization reduces excessive dependence
on individual feature channels in deeper semantic representations.

\subsection{\texorpdfstring{ Lightweight residual
bottleneck}{ Lightweight residual bottleneck}}\label{lightweight-residual-bottleneck}

The deepest encoder representation F4 is first processed by a
lightweight residual bottleneck. A 1×1 convolution reduces its channel
dimension from 768 to 192. A depthwise separable 3×3 convolution
performs spatial transformation in the reduced feature space, after
which another 1×1 convolution restores the channel dimension to 768.

\emph{\(B_0 = F_4 + \gamma_b H(F_4).\)}

The transformation H(·) includes pointwise channel reduction, Layer
Normalization, GELU activation, depthwise separable convolution,
dropout, pointwise channel restoration, and final Layer Normalization.
The learnable scalar \(\gamma_b\) is initialized to zero, so the bottleneck
initially behaves as an identity mapping and progressively introduces
its residual transformation during optimization.

\subsection{\texorpdfstring{ Multi-Scale Context
module}{ Multi-Scale Context module}}\label{multi-scale-context-module}

The Multi-Scale Context module is applied to B0 to enrich the bottleneck
representation using multiple receptive fields while maintaining
moderate computational complexity. A 1×1 convolution first reduces the
768-channel input to 96 channels.

\emph{R = GELU(LN(Conv1×1(B0))).}

Three parallel depthwise 3×3 convolutions with dilation rates \(d \in \{1,
3, 5\}\) are applied:

\emph{\(C_d = \operatorname{GELU}(\operatorname{LN}(\operatorname{DWConv}_{3\times3,d}(R))),\ d \in \{1, 3, 5\}.\)}

A fourth branch extracts image-level context through global average
pooling:

\emph{Cg = Up(GELU(Conv1×1(GAP(R)))).}

The four branches are concatenated and projected back to 768 channels:

\emph{C = {[}C1, C3, C5, Cg{]},}

\emph{\(B = B_0 + \gamma_c P(C).\)}

The projection P(·) consists of a 1×1 convolution, Layer Normalization,
GELU activation, and dropout. The scalar \(\gamma_c\) is initialized to zero.
Dilation rates 1, 3, and 5 provide local and progressively broader
neighborhoods, while the pooled branch supplies image-level context.

\subsection{Boundary-Sensitive Enhancement and Integration
module}\label{boundary-sensitive-enhancement-and-integration-module}

At each skip-connected decoder stage, the upsampled decoder feature is
concatenated with the corresponding encoder feature. The resulting
representation is refined using BSEI. Let Sl = {[}Ul, Fl{]} denote the
concatenated decoder and encoder features at level l.

\emph{Ql = GELU(LN(Conv1×1(Sl))),}

\emph{Zl = GELU(LN(SepConv3×3(Ql))).}

A learned depthwise 3×3 convolution generates a local response map. The
absolute value represents the magnitude of the learned spatial
variation, and a sigmoid bounds the response:

\emph{Al = sigmoid(\textbar DWConv3×3(Zl)\textbar).}

The enhanced feature is calculated by residual modulation:

\emph{\(Y_l = Z_l + A_l \odot Z_l.\)}

Locations with stronger learned spatial responses therefore receive
greater amplification. This operation should be described as learned
boundary-sensitive or edge-response enhancement, rather than explicit
boundary detection. It does not use Sobel filters, manually generated
edge maps, channel squeeze-and-excitation, or boundary annotations.

\subsection{Progressive decoder}\label{progressive-decoder}

Each decoder block consists of bilinear up-sampling by a factor of two,
a depthwise separable 3×3 convolution, Layer Normalization, GELU
activation, and spatial dropout. The first decoder stage transforms the
context-enhanced bottleneck from 768 to 384 channels. After
concatenation with F3, BSEI produces a 384-channel representation.

\emph{D4 = M\_BSEI,4({[}M\_dec,4(B), F3{]}),}

\emph{D3 = M\_BSEI,3({[}M\_dec,3(D4), F2{]}),}

\emph{D2 = M\_BSEI,2({[}M\_dec,2(D3), F1{]}),}

\emph{D1 = M\_BSEI,1(M\_dec,1(D2)).}

The output channel dimensions of D3 and D2 are 192 and 96, respectively.
The final decoder block increases the spatial resolution from H/4×W/4 to
H/2×W/2. No encoder skip feature is concatenated at this final stage
because the encoder stem begins at one-quarter input resolution.

\subsection{High-resolution Detail
Branch}\label{high-resolution-detail-branch}

The Detail Branch processes the original input independently of the deep
encoder. Its objective is to retain relatively high-resolution
appearance information that may be attenuated by the stride-4 encoder
stem. The branch begins with a 3×3 convolution with stride 2, producing
a 32-channel feature at one-half input resolution.

\emph{T0 = GELU(BN(Conv3×3,s=2(I))),}

\emph{T = GELU(BN(SepConv3×3(T0))).}

Dropout is applied after both operations. Unlike the main encoder and
decoder, this branch uses Batch Normalization.

\subsection{Gated Detail Fusion}\label{gated-detail-fusion}

The final decoder feature D1 and the detail representation T have the
same spatial resolution. Before fusion, the detail feature is projected
using a depthwise separable convolution:

\emph{\(\hat{T} = \operatorname{GELU}(\operatorname{BN}(\operatorname{SepConv}_{3\times3}(T))).\)}

A spatially varying gate is estimated from the concatenated decoder and
detail representations:

\emph{\(G = \operatorname{sigmoid}(W_{g2}(\operatorname{GELU}(W_{g1}([D_1, \hat{T}])))).\)}

The gate contains 32 channels and modulates the detail representation:

\emph{\(T_g = G \odot \hat{T}.\)}

The gated detail representation is concatenated with the decoder feature
and projected to 128 channels:

\emph{Ffuse = GELU(LN(Conv1×1({[}D1, Tg{]}))).}

The gate controls the spatial and channel-wise contribution of the
shallow detail representation before it enters the final prediction
head.

\subsection{Prediction head and deep
supervision}\label{prediction-head-and-deep-supervision}

The fused feature is processed by two depthwise separable convolutions.
The first reduces the channel dimension from 128 to 96, and the second
reduces it from 96 to 48. A 1×1 convolution then produces the
segmentation logits.

\emph{P0 = Conv1×1(R(Ffuse)),}

\emph{P = UpH,W(P0).}

For binary segmentation, the probability map at inference is obtained as
\(\hat{Y} = \operatorname{sigmoid}(P)\). The implementation returns logits rather than sigmoid
probabilities; consequently, the activation should be applied by the
loss function or by the inference procedure.

During training, auxiliary 1×1 prediction heads are attached to D2, D3,
and D4. Their logits are resized to the original input resolution:

\emph{P2 = UpH,W(Conv1×1(D2)),}

\emph{P3 = UpH,W(Conv1×1(D3)),}

\emph{P4 = UpH,W(Conv1×1(D4)).}

A general deeply supervised objective can be written as:

\emph{Ltotal = Lseg(P,Y) + $\lambda$2 Lseg(P2,Y) + $\lambda$3 Lseg(P3,Y) + $\lambda$4
Lseg(P4,Y).}

\emph{Here, Y denotes the reference mask and $\lambda$2, $\lambda$3, and $\lambda$4 control the
auxiliary losses. The exact segmentation loss and numerical auxiliary
weights are not defined in the supplied architecture file and must be
obtained from the training implementation before the manuscript is
finalized.}

\subsection{Model complexity}\label{model-complexity}

For a single-class output configuration, the supplied implementation
contains 29,609,302 trainable parameters (approximately 29.61 million).
This value is consistent with the approximately 29.60 million parameters
reported in the later complexity table of the experimental draft. Any
earlier table reporting 38.5 million parameters should be corrected or
explicitly identified as belonging to a different model version.

